\documentclass[11pt]{article}
\usepackage[T1]{fontenc}
\usepackage{lmodern,amsmath,amssymb,amsthm,microtype,booktabs}
\usepackage[a4paper,margin=26mm]{geometry}
\usepackage[colorlinks=true,urlcolor=blue,citecolor=blue]{hyperref}
\newtheorem{theorem}{Theorem}
\newtheorem{lemma}[theorem]{Lemma}
\newtheorem{corollary}[theorem]{Corollary}
\newcommand{\E}{\mathbb E}
\newcommand{\R}{\mathbb R}
\newcommand{\TV}{\mathrm{TV}}
\newcommand{\norm}[1]{\left\|#1\right\|_1}
\title{The Sharp Logarithmic Coefficient\protect\\of a Controlled Information Envelope}
\author{Jan Mikul\'ik}
\date{12 October 2026\quad---\quad Research manuscript}
\begin{document}
\maketitle
\begin{abstract}
For finitely many centered, finitely supported increment laws with a
common positive variance, we identify the logarithmic displacement of
their adaptive lower CDF envelope. Its coefficient is the range of the
third moments divided by $3v\sqrt{2\pi e}$. The proof uses the established
logarithmic envelope-width bound, an explicit Ornstein--Uhlenbeck
corrector, and an integrated stability estimate. It does not differentiate
the actual Bellman envelope. The resulting error is
$O(\sqrt{\log n\log\log n})$, both in the mean and in unscaled
$W_1$ distance from a translated Gaussian. The finite causal entropy
sandwich transfers the coefficient to optimal seed entropy.
\end{abstract}

\paragraph{Status and dependency.}
This is an analytic proof manuscript, not a Lean certificate or an
independently reviewed paper. It uses the logarithmic moment bound from
the accompanying repository as an explicit input, stated below.
No priority claim is made. In particular, an $O(1)$ remainder is
\emph{not} proved here.

\section{Statement}
Let $\mu_1,\ldots,\mu_d$ be finitely supported probability laws on $\R$,
with
\[
 \E X_i=0,\qquad \E X_i^2=v>0,\qquad
 \kappa_i=\E X_i^3.
\]
Put $\kappa_- =\min_i\kappa_i$, $\kappa_+=\max_i\kappa_i$ and
\[
 A_i f(x)=\E f(x-X_i),\quad Tf=\min_i A_i f,\quad
 Uf=\max_i A_i f,\quad H(x)=\mathbf1_{[0,\infty)}(x).
\]
The $n$-step lower and upper policy CDF envelopes are
$G_n=T^nH$ and $V_n=U^nH$. Policies may depend on the complete observed
history. Their terminal increment sums have CDFs between $G_n$ and $V_n$.
Write $\nu_n$ for the law with CDF $G_n$, and
$e_n=\int_\R(H-G_n)$ for its mean. Let $\phi,\Phi$ denote the standard
normal density and CDF. All logarithms in asymptotic terms are natural.

\begin{theorem}\label{thm:main}
For every fixed family as above, with
\begin{equation}\label{eq:c}
 c=\frac{\kappa_+-\kappa_-}{3v\sqrt{2\pi e}},
\end{equation}
one has, as $n\to\infty$,
\begin{align}
 W_1\bigl(\nu_n,N(c\ln n,vn)\bigr)
   &=O\bigl(\sqrt{\ln n\ln\ln n}\bigr),\label{eq:W1}\\
 e_n&=c\ln n+O\bigl(\sqrt{\ln n\ln\ln n}\bigr).
 \label{eq:mean}
\end{align}
The constants may depend on the fixed increment family. No lattice,
aperiodicity or nonlattice hypothesis is imposed.
\end{theorem}

\begin{corollary}[Causal seed entropy]\label{cor:cost}
For a finite stationary kernel whose action laws $P_i$ have common
entropy $h$ and common positive information variance $v$, set
$\kappa_i=\E(-\log_2P_i(Y_i)-h)^3$. In bits,
\[
 \boxed{C_n=nh+
 \frac{\kappa_+-\kappa_-}{3v\sqrt{2\pi e}}\ln n
 +O\bigl(\sqrt{\ln n\ln\ln n}\bigr).}
\]
The same expansion holds for the lower envelope mean and for any exact
seeds whose entropy is within a bounded additive amount of that mean.
\end{corollary}
\begin{proof}
Center each action's information by $h$. The centered envelope is $G_n$,
and $L_n-nh=e_n$. Apply Theorem~\ref{thm:main} and the finite-tree
sandwich $L_n\le C_n\le L_n+c_\star$ from~\cite{repo}.
\end{proof}

\section{The existing input and integrated stability}
\paragraph{Input: logarithmic moment bound.}
For every finite centered increment family of common positive variance,
the repository proves $0\le e_n\le K(1+\ln(n+1))$.
The declaration is
\texttt{CausalSpectrum.Sqrt.IncFamily.eG\_log\_upper} in
\texttt{MomentAssembly.lean}. Its nondegeneracy assumption is automatic:
a finite law of positive variance has two distinct atoms, a positive
equal-mass part of which is a translated symmetric two-point law.
This input is an order bound, not an assumption of the coefficient to be proved.

Apply the same result to the reflected increment family $(-X_i)_i$.
Reflection exchanges the upper and lower envelopes, apart from irrelevant
values at their finitely many jump points. Consequently
\begin{equation}\label{eq:width}
 w_n:=\int_\R(V_n-G_n)=e_n+e_n^{\mathrm{ref}}
 \le K(1+\ln(n+1)).
\end{equation}
In particular $\bar w_n:=\max_{0\le j\le n}w_j=O(\ln(n+2))$.

\begin{lemma}[Integrated amplification]\label{lem:amp}
For $r\in L^1(\R)\cap BV(\R)$,
\begin{equation}\label{eq:amp}
 \int_\R U^n|r|(x)\,dx\le\norm{r}+w_n\TV(r).
\end{equation}
For bounded functions $f,g$ with $f-g\in L^1\cap BV$,
\begin{equation}\label{eq:stability}
 \norm{T^nf-T^ng}\le\norm{f-g}+w_n\TV(f-g).
\end{equation}
\end{lemma}
\begin{proof}
Fix a reference policy repeating one action, and let $F_0$ be its
$n$-step CDF. For any policy $\pi$, its CDF $F_\pi$ and $F_0$ lie
between $G_n$ and $V_n$. For absolutely continuous $r$, integration by
parts, with $r(\pm\infty)=0$, gives
\[
 \E_\pi|r(x-S_n)|\le
 \E_0|r(x-S_n)|+
 \int_\R |r'(x-y)|(V_n(y)-G_n(y))\,dy.
\]
The bound is independent of $\pi$. Take the supremum over policies,
integrate in $x$, and apply Tonelli. This proves~\eqref{eq:amp} for
absolutely continuous $r$. Mollification proves the $BV$ version:
variation does not increase; for a fixed horizon there are finitely
many policies, so their finite-translation operators are continuous on
$L^1$, and the supremum passes to the limit in $L^1$.
Finally $|T^nf-T^ng|\le U^n|f-g|$ pointwise, which proves
\eqref{eq:stability}.
\end{proof}

\section{An explicit resonant corrector}
Define
\begin{align}
 f(z)&=\min_i\left[-\frac{\kappa_i}{6v^{3/2}}
                      (z^2-1)\phi(z)\right],\label{eq:forcing}\\
 a&=\int_\R f(z)\,dz=-\frac{\kappa_+-\kappa_-}{3v^{3/2}}\phi(1)
       =-\frac{c}{\sqrt v},\label{eq:a}\\
 g(z)&=2(a\phi(z)-f(z)),\qquad
 J(z)=\int_{-\infty}^z g(y)\,dy,\label{eq:J}\\
 B(z)&=\phi(z)\int_0^z\frac{J(y)}{\phi(y)}\,dy.\label{eq:B}
\end{align}
Indeed $\int(z^2-1)\phi=0$ and
$\int|z^2-1|\phi=4\phi(1)$, proving~\eqref{eq:a}.
Since $\int g=0$, formula~\eqref{eq:B} decays at both ends and solves
\begin{equation}\label{eq:ODE}
 B''+zB'+B=2(a\phi-f).
\end{equation}

\paragraph{Regularity and tails, explicitly.}
The function $f$ is continuous, smooth off $\{-1,1\}$, and has only
derivative jumps there. Thus $B\in C^2$, with locally absolutely
continuous second derivative. For $z<-1$ and $z>1$, respectively,
\[
 \frac{J(z)}{\phi(z)}=
 2a\frac{\Phi(z)}{\phi(z)}-\frac{\kappa_+z}{3v^{3/2}},
 \qquad
 \frac{J(z)}{\phi(z)}=
 -2a\frac{1-\Phi(z)}{\phi(z)}-\frac{\kappa_+z}{3v^{3/2}}.
\]
The elementary Gaussian tail inequality
$1-\Phi(z)\le\phi(z)/z$ for $z>0$, and its reflected version,
show that the integral in~\eqref{eq:B} is $O(1+z^2)$.
The equation $B'+zB=J$ and~\eqref{eq:ODE}, followed by its
almost-everywhere derivative, then show that $B,B',B'',B'''$ are
bounded by a polynomial times $\phi$ on each tail. All their
polynomially weighted absolute integrals are finite. In particular
$B\in W^{3,1}(\R)$. No fourth classical derivative of $B$ is required.

For $t\ge1$, put
\begin{equation}\label{eq:profile}
 z=\frac{x-c\ln t}{\sqrt{vt}},\qquad
 p_t(x)=\Phi(z),\quad q_t(x)=t^{-1/2}B(z),\quad u_t=p_t+q_t.
\end{equation}
The function $u_t$ need not be a CDF; it is used only as a bounded
comparison profile. Its limiting values are zero and one. Direct
differentiation using~\eqref{eq:ODE} gives
\begin{equation}\label{eq:PDE}
 \partial_tu_t-\frac v2\partial_x^2u_t
 =t^{-3/2}f(z)-\frac{c}{\sqrt v}\,t^{-2}B'(z).
\end{equation}
The translation $c\ln t$ absorbs the resonant Gaussian mode exactly.

\section{Summable discrete residuals}
\begin{lemma}\label{lem:res}
There is a finite $K$, depending only on the increment family, such that
for every real $t\ge1$, the residual $R_t=u_{t+1}-Tu_t$ belongs to
$W^{1,1}(\R)$ and
\begin{equation}\label{eq:res}
 \norm{R_t}+\norm{R_t'}\le Kt^{-3/2}.
\end{equation}
\end{lemma}
\begin{proof}
We give the derivative estimates as well as the value estimate, since
the derivative is essential to integrated stability. The regularity
above and the change of variables $dx=\sqrt{vt}\,dz$ imply
\begin{align*}
 \norm{\partial_x^k p_t}&=O(t^{-(k-1)/2})\quad(k\ge1),\\
 \norm{\partial_x^k q_t}&=O(t^{-k/2})\quad(0\le k\le3),\\
 \norm{\partial_t^2u_t}&=O(t^{-3/2}),\qquad
 \norm{\partial_t\partial_xq_t}=O(t^{-3/2}).
\end{align*}
For the last two estimates, use
$\partial_tz=-z/(2t)-c/(\sqrt v\,t^{3/2})$; only weighted
$B,B',B''$ occur. Also
$\partial_tp_t=(v/2)p_t''-(c/t)p_t'$.
All $O$ constants here are finite and uniform for $t\ge1$.

Taylor's formula under translation, in $L^1$, gives for each action
\[
 A_ip_t=p_t+\frac v2p_t''-\frac{\kappa_i}{6}p_t'''+e_{i,t},
 \qquad \norm{e_{i,t}}\le\frac{\E|X_i|^4}{24}\norm{p_t''''}
 =O(t^{-3/2}),
\]
and
\[
 A_iq_t=q_t+\frac v2q_t''+b_{i,t},\qquad
 \norm{b_{i,t}}\le\frac{\E|X_i|^3}{6}\norm{q_t'''}
 =O(t^{-3/2}).
\]
The scalar minimum is Lipschitz with respect to the maximum norm.
Since $\min_i[-\kappa_i p_t'''/6]=t^{-3/2}f(z)$, summing these
error bounds yields
\[
 \norm{Tu_t-u_t-\tfrac v2u_t''-t^{-3/2}f(z)}=O(t^{-3/2}).
\]
Taylor expansion in time,~\eqref{eq:PDE}, and
$\norm{t^{-2}B'(z)}=O(t^{-3/2})$ prove the estimate for $R_t$.

For its derivative, write
$d_{i,t}=A_i u_t-u_t-(v/2)u_t''$.
Taylor expansion to order two for $p_t'$ gives
\[
 \norm{A_ip_t'-p_t'-\tfrac v2p_t'''}
 \le\frac{\E|X_i|^3}{6}\norm{p_t''''}=O(t^{-3/2}).
\]
Centered Taylor expansion to order one for $q_t'$, followed by
the triangle inequality, gives
\[
 \norm{A_iq_t'-q_t'-\tfrac v2q_t'''}
 \le v\norm{q_t'''}=O(t^{-3/2}).
\]
Hence $\norm{d_{i,t}'}=O(t^{-3/2})$.
The minimum of finitely many $W^{1,1}$ functions has derivative
bounded almost everywhere by the sum of the absolute derivatives.
Finally
\[
 \norm{p_{t+1}'-p_t'-\tfrac v2p_t'''}=O(t^{-3/2}).
\]
To verify this last estimate, integrate
$\partial_s p_s'=(v/2)p_s'''-(c/s)p_s''$ on $[t,t+1]$;
$\norm{p_s''}=O(s^{-1/2})$ and
$\norm{\partial_s p_s'''}=O(s^{-2})$.
For $q_t$, the preceding time-derivative bound and
$\norm{q_t'''}=O(t^{-3/2})$ give the same estimate.
Differentiating
$R_t=u_{t+1}-u_t-(v/2)u_t''-\min_i d_{i,t}$ now proves
the derivative bound. All derivatives in this paragraph may be read
as weak derivatives, so the two switching points cause no missing term.
\end{proof}

\section{From a late-start profile to the actual envelope}
For any $t_0\ge1$, successive one-step comparisons and positivity of $U$
give the pointwise bound
\[
 |T^nu_{t_0}-u_{t_0+n}|
 \le\sum_{j=0}^{n-1}U^{n-1-j}|R_{t_0+j}|.
\]
Here one uses $|Tf-Tg|\le U|f-g|$ and subadditivity of $U$.
By Lemmas~\ref{lem:amp} and~\ref{lem:res},
\begin{equation}\label{eq:prop}
 \norm{T^nu_{t_0}-u_{t_0+n}}
 \le K\frac{1+\bar w_n}{\sqrt{t_0}}.
\end{equation}
Thus a summable residual is sufficient after its derivative's
amplification by the envelope width has also been accounted for.

It remains to compare the step initial condition with a Gaussian.
Let $s=\sqrt{vt_0}$, $b(x)=\Phi(x/s)$, choose $R>0$, and put
$A=sR$, $\varepsilon=1-\Phi(R)$. The pointwise inequalities
\[
 H(x)\le b(x+A)+r_+(x),\qquad
 b(x-A)\le H(x)+r_-(x)
\]
hold with
$r_+(x)=\mathbf1_{x\ge0}(1-b(x+A))$ and
$r_-(x)=\mathbf1_{x<0}b(x-A)$.
They satisfy
\[
 \norm{r_\pm}\le s\phi(R),\qquad \TV(r_\pm)=2\varepsilon.
\]
Propagate the inequalities with $T$, use translation invariance and
Lemma~\ref{lem:amp}, and integrate the positive and negative parts.
Since $T^nb$ is a CDF, its integral difference from a translate by $A$
is exactly $A$. Therefore
\begin{equation}\label{eq:initial}
 \norm{G_n-T^nb}\le2A+2s\phi(R)+4\varepsilon w_n.
\end{equation}
The Gaussian $p_{t_0}$ is $b$ translated by $c\ln t_0$, so
$\norm{T^np_{t_0}-T^nb}=|c\ln t_0|$.
Furthermore $\norm{q_{t_0}}=\sqrt v\norm{B}$ and
$\TV(q_{t_0})=t_0^{-1/2}\norm{B'}$. Thus
\begin{equation}\label{eq:initialB}
 \norm{T^nu_{t_0}-T^np_{t_0}}
 \le \sqrt v\norm B+w_n t_0^{-1/2}\norm{B'}.
\end{equation}

Take $t_0=\ln(n+e^2)$ and $R=\sqrt{2\ln t_0}$.
Then $\phi(R)=(2\pi)^{-1/2}/t_0$,
$\varepsilon\le\phi(R)/R$, and $\bar w_n=O(t_0)$ by
\eqref{eq:width}. Equations~\eqref{eq:prop}--\eqref{eq:initialB} imply
\begin{equation}\label{eq:last}
 \norm{G_n-u_{n+t_0}}
 =O\bigl(\sqrt{\ln n\ln\ln n}\bigr).
\end{equation}
Since $\norm{u_t-p_t}=\sqrt v\norm B$, this also holds with
$p_{n+t_0}$ in place of $u_{n+t_0}$.
Coupling Gaussians with the same standard normal variable shows
\[
 W_1\bigl(N(c\ln(n+t_0),v(n+t_0)),N(c\ln n,vn)\bigr)
 \le |c|\ln(1+t_0/n)
 +\sqrt{2v/\pi}\,(\sqrt{n+t_0}-\sqrt n)=o(1).
\]
The one-dimensional identity $W_1(\mu,\nu)=\int|F_\mu-F_\nu|$
proves~\eqref{eq:W1}. The difference of means is at most $W_1$,
proving~\eqref{eq:mean} and completing Theorem~\ref{thm:main}.\hfill$\square$

\section{A concrete dyadic information pair}
Take
\[
 P=(1/4,\underbrace{1/16,\ldots,1/16}_{12}),\qquad
 Q=(\underbrace{1/8,\ldots,1/8}_{6},
       \underbrace{1/32,\ldots,1/32}_{8}).
\]
Their entropies are $h=7/2$, their information variances are $v=3/4$,
and their centered third moments are $-3/4$ and $3/4$.
Every atom is a power of $1/2$, so dyadic exactness gives $C_n=L_n$
at every horizon. Hence
\[
 \boxed{C_n=\frac72n+
 \frac{2}{3\sqrt{2\pi e}}\ln n
 +O\bigl(\sqrt{\ln n\ln\ln n}\bigr).}
\]
The coefficient is $0.161313816346\ldots$ bits per unit of $\ln n$.
The result is about the unrestricted optimal causal seed, not a
particular block construction.

\paragraph{Scope.}
Only the extreme third moments enter the logarithmic coefficient.
Intermediate actions and higher moments can affect the remainder.
When the third moments all agree, the existing moment classification
already gives the stronger bounded-excess conclusion; this manuscript
is principally useful when $\kappa_+>\kappa_-$.
It does not prove a bounded remainder, a pointwise $1/n$ limit for
$e_{n+1}-e_n$, or a uniform theorem as $v\downarrow0$ or the kernel varies.
The profiles $u_t$ are auxiliary functions, not asserted probability laws.

\paragraph{Related work and provenance.}
Rates in martingale central limit theorems and nonlinear central limit
theorems have a substantial literature; see~\cite{mourrat,fang}.
Those references provide context, not the coefficient calculation in
this proof. The nonstandard analytic input here is specifically the
\emph{integrated envelope-width} bound~\eqref{eq:width}, supplied by
the repository's moment theorem. A bound for each fixed policy in
$W_1$ alone would not justify taking the pointwise envelope.

\begin{thebibliography}{9}
\bibitem{repo} J.~Mikul\'ik, \emph{Finite Shannon theorem},
\url{https://github.com/jendamikulik/finite-Shannon-theorem},
retrieved 12 October 2026. Analytic input:
\texttt{CausalSeed/MomentAssembly.lean}, theorem
\texttt{Sqrt.IncFamily.eG\_log\_upper}; compare also
\texttt{MomentWidth.lean}, theorem \texttt{polE\_abs\_le}.
The new corrector and coefficient theorem in this manuscript are not
claimed to be formalized there.
\bibitem{mourrat} J.-C.~Mourrat, On the rate of convergence in the
martingale central limit theorem, \emph{Bernoulli} \textbf{19}(2)
(2013), 633--645. \href{https://doi.org/10.3150/12-BEJ417}{doi:10.3150/12-BEJ417}.
\bibitem{fang} X.~Fang, S.~Peng, Q.-M.~Shao and Y.~Song,
\emph{Limit theorems with rate of convergence under sublinear expectations},
\url{https://arxiv.org/abs/1711.10649}.
\end{thebibliography}
\end{document}
